\section{MC modeling, including polarisation (Day 2, morning)}

\subsection{Measuring the polarization of boosted, hadronic $W$ bosons with jet substructure observables}
\speaker{Songshaptak De}{Jo\v{z}ef Stefan Institute; with V.~Rentala and W.~Shepherd, JHEP 05 (2025) 028}
A new, production-process-independent technique to extract the longitudinal and
transverse polarization fractions of boosted hadronic $W$ bosons was presented.
The key object is a jet-substructure observable $p_\theta =
\Delta E^{\mathrm{reco}}/p_W^{\mathrm{reco}}$, the subjet energy difference
normalized to the fat-jet momentum, which at parton level equals $|\costh|$,
the decay polar angle in the $W$ rest frame; it is shown to have smaller
reconstruction errors than the traditionally used opening-angle proxy,
especially at large boosts, and remains faithful up to $|\costh|\lesssim0.9$.
Hadron-level 2D templates in $(p_\theta,\tau_{21})$ are built for longitudinal
and transverse $W$ jets (scalar/pseudoscalar $\phi\to WW\to 4q$ samples,
$800<\pT<1000$~GeV, soft-drop groomed anti-$k_t$ $R=1.0$ jets) and for the QCD
dijet background, and fitted to pseudo-data. Applied to semi-leptonic $WW$ VBS
at the HL-LHC, the method reconstructs the \emph{transverse} fraction to
20--30\% with 3--10~ab$^{-1}$, but with SM production alone the
\emph{longitudinal} fraction cannot be distinguished from zero even with
10~ab$^{-1}$ ($f_L = 0.0075 \pm 0.0045$ for a true 0.0074), because of the
$\mathcal{O}(25{:}1)$ QCD background and the tiny $LL$ rate. The technique
would however be sensitive to BSM-enhanced longitudinal production, and can be
strengthened by using $p_\theta$ inside a multivariate analysis.

\subsection{A systematic approach to amplitude-assisted polarization tagging (remote)}
\speaker{Erik Bachmann}{TU Dresden; with D.~M.~Hoppe}
This theory contribution addressed the question of what the \emph{optimal}
polarization observable is. In contrast to classification-based strategies,
which approximate polarization fractions through relative frequencies of purely
polarized training samples, amplitude-assisted polarization tagging learns
per-event polarization fractions directly from the matrix elements. The talk
presented a systematic construction of the corresponding optimal observable,
which defines the fundamental limit for separating polarization components and
therefore provides a natural benchmark against which neural-network approaches
can be tested---enabling a diagnosis of where and why their performance falls
short of optimality. The methodology was illustrated on the simplest
non-trivial process, leptonic $Z$+jet production at leading order, and
extensions to $W$ and diboson processes as well as to hadronic decay channels
were discussed.

\subsection{MC roundtable and discussion}
\speaker{Emanuele Re, Giovanni Pelliccioli}{Universit\`a \& INFN Milano-Bicocca, conveners}
The roundtable, continuing the thread opened by the Day-1 theory review,
discussed the status and interoperability of polarized event generation:
polarized samples in MadGraph (LO and multi-jet merged), SHERPA (nLO+PS),
POWHEG-BOX-RES (NLO+PS, and now SMEFT@NLO+PS), and fixed-order DPA computations
(MoCaNLO, NNLO~QCD frameworks). Central topics were the definition dependence
of polarized signals (frame choices, NWA vs DPA vs off-shell, interference and
Breit--Wigner treatments), the size of higher-order corrections to hadronic
decays and their interplay with parton-shower matching in the boosted regime,
and the recommendation that experiments quote results with clearly documented
polarization conventions to allow reinterpretation. The ongoing COMETA
comparative study of polarized $ZZ$ (and its planned extensions to $WZ$ and to
semi-leptonic final states) was identified as the natural vehicle to converge
on best practices and best predictions.
